% Unit 128 terjemahan Bahasa Indonesia dari chapter9.tex baris 1257--1367.
% Sumber: Methods of Algebra, Volume 2, commit
% 9a5803ff77dd3257484cb177f851a73770a59dd3 (CC BY 4.0).
% stable-unit-id: o014.aljabr2.chapter9.reconstruction-a
% Cakupan: Bagian 9.5 hingga Lema rekonstruksi dengan generator projektif.
% Koreksi sumber terpaksa: baris 1275 memakai kategori modul kanan
% cated{Mod}A; baris 1331 mensyaratkan N'_2 memuat Image(f|N'_1), bukan
% Image omega(f|N'_1), yang bertipe sebagai submodul omega(N'_2).

% segment-id: o014.li2.chapter9.reconstruction-a.h001
\section{Teorema rekonstruksi}\label{sec:Tannaka-coend}
\hypertarget{o014.aljabr2.chapter9.reconstruction-a}{ }
% segment-id: o014.li2.chapter9.reconstruction-a.p001
Sepanjang bagian ini, $\Bbbk$ adalah gelanggang komutatif tetap. Kecuali
dinyatakan lain, semua kategori dan funktor yang dibahas di bawah bersifat
$\Bbbk$-linear; khususnya, bifunktor $\otimes$ dari kategori monoidal
dengan sendirinya $\Bbbk$-linear dalam setiap variabel. Untuk aljabar-$\Bbbk$
$B$, Konvensi~\ref{con:Mod-pfg} telah mendefinisikan kategori kecil secara
esensial
$\catesubd{Mod}{\mathrm{pfg}}B\subset
\catesubd{Mod}{\mathrm{fg}}B$. Jika $C$ koaljabar dalam kategori monoidal
$(B,B)\dcate{Mod}$, tuliskan $\cated{Comod}C$ untuk kategori komodul-$C$
kanan; struktur komodul ini ditumpangkan pada struktur dalam
$\cated{Mod}B$, seperti dijelaskan setelah Definisi~\ref{def:cogebroid}.
Dengan memberlakukan syarat projektivitas dan keterhinggaan sebagai
modul-$B$ kanan, diperoleh subkategori penuh
% segment-id: o014.li2.chapter9.reconstruction-a.q001
\[ \underbracket{\catesubd{Comod}{\mathrm{pf}}C}_{\text{projektif + terbangkit hingga}} \subset \underbracket{\catesubd{Comod}{\mathrm{f}}C}_{\text{terbangkit hingga}} \subset \cated{Comod}C. \]
\index[sym1]{Comodpfg@$\cate{Comod}_{\mathrm{pf}}, \cate{Comod}_{\mathrm{f}}$}

% segment-id: o014.li2.chapter9.reconstruction-a.p002
Syarat terbangkit hingga memastikan bahwa
$\catesubd{Comod}{\mathrm f}C$ juga kategori kecil secara esensial dalam
arti Definisi~\ref{def:essentially-small-cat}.

% segment-id: o014.li2.chapter9.reconstruction-a.p003
Untuk memperoleh pemahaman lebih konkret tentang
$\coEnd(\omega)=\coEnd_{\Bbbk}(\omega)$ dari \S\ref{sec:coend}, sekarang
tinjau situasi Proposisi~\ref{prop:Morita-comonad}. Ambil
aljabar-$\Bbbk$ $A,B$, yang tidak harus komutatif, dan bimodul-$(A,B)$
$P$. Andaikan $P$ projektif terbangkit hingga sebagai modul-$B$ kanan.
Ambil $\mathcal{A}$ sebagai subkategori penuh $\cated{Mod}A$ dengan
syarat
\begin{itemize}
	\item $A\in\Obj(\mathcal{A})$;
	\item $\omega:=(\mathord\cdot)\dotimes{A}P$ memberikan funktor
	$\mathcal{A}\to\catesubd{Mod}{\mathrm{pfg}}B$.
\end{itemize}

% segment-id: o014.li2.chapter9.reconstruction-a.p004
Proposisi~\ref{prop:Morita-comonad} menyatakan bahwa bimodul-$(B,B)$
$P^\vee\dotimes{A}P$ adalah koaljabar yang bersesuaian dengan komonad yang
ditentukan oleh pasangan adjoin
% segment-id: o014.li2.chapter9.reconstruction-a.q002
\begin{equation*}\begin{tikzcd}
	(\cdot) \dotimes{A} P : \cated{Mod}A \arrow[shift left, r] & \cated{Mod}B: (\cdot) \dotimes{B} P^\vee \arrow[shift left, l].
\end{tikzcd}\end{equation*}
Koaljabar $P^\vee\dotimes{A}P$ secara alami berkoaksi pada setiap
$\omega(N)=N\dotimes{A}P$. Secara eksplisit, koaksinya ialah
% segment-id: o014.li2.chapter9.reconstruction-a.q003
\[ N \dotimes{A} P \xrightarrow{\identity_N \otimes \mathrm{coev} \otimes \identity_P} N \dotimes{A} P \dotimes{B} P^\vee \dotimes{A} P, \quad N \in \Obj(\cated{Mod}A). \]

% segment-id: o014.li2.chapter9.reconstruction-a.p005
Definisi--Proposisi~\ref{def:cogebra-L} kemudian memberikan homomorfisme
koaljabar kanonik
% segment-id: o014.li2.chapter9.reconstruction-a.q004
\[ \coEnd(\omega) \to P^\vee \dotimes{A} P. \]

% segment-id: o014.li2.chapter9.reconstruction-a.pr001
\begin{proposition}\label{prop:endomorphism-vs-P}
	Dalam situasi di atas, homomorfisme
	$\coEnd(\omega)\to P^\vee\dotimes{A}P$ adalah isomorfisme.
\end{proposition}
\begin{proof}
	Ambil subkategori penuh $\mathcal{A}^\flat$ dari $\cated{Mod}A$ dengan
	$\Obj(\mathcal{A}^\flat)=\{A\}$, dan tuliskan $i$ untuk funktor
	inklusinya. Catatan~\ref{rem:L-cogebra-functoriality} memberikan diagram
	komutatif
	% segment-id: o014.li2.chapter9.reconstruction-a.q005
	\[\begin{tikzcd}[column sep=tiny]
		\coEnd(\omega i) \arrow[rr] \arrow[rd] & & \coEnd(\omega), \arrow[ld] \\
		& P^\vee \dotimes{A} P &
	\end{tikzcd}\]
	Pembaca dapat memeriksa langsung bahwa
	$\coEnd(\omega i)\rightiso P^\vee\dotimes{A}P$. Jadi
	$\coEnd(\omega i)\to\coEnd(\omega)$ memiliki invers kiri. Jika dapat
	ditunjukkan bahwa morfisme ini surjektif, maka
	$\coEnd(\omega i)\rightiso\coEnd(\omega)$ dan akibatnya
	$\coEnd(\omega)\rightiso P^\vee\dotimes{A}P$.

	Menurut konstruksi konkret $\coEnd(\omega)$ dan $\coEnd(\omega i)$,
	masalahnya ialah membuktikan, untuk setiap $Y\in\Obj(\mathcal{A})$,
	% segment-id: o014.li2.chapter9.reconstruction-a.q006
	\[ \Image\left[ \omega(Y)^\vee \dotimes{\Bbbk} \omega(Y) \to \coEnd(\omega) \right] \subset \Image\left[ \omega(A)^\vee \dotimes{\Bbbk} \omega(A) \to \coEnd(\omega) \right]. \]

	Tuliskan $t\in\omega(Y)^\vee\dotimes{\Bbbk}\omega(Y)$ sebagai jumlah
	hingga $\sum_i\lambda_i\otimes o_i$. Karena bentuk konkret $\omega$,
	setiap $o_i\in\omega(Y)$ adalah kombinasi linear hingga unsur berbentuk
	$\omega(f)(x)$, dengan $x\in\omega(A)$ dan
	$f\in\Hom_{\mathcal{A}}(A,Y)$. Jadi masalah direduksi lebih lanjut ke
	kasus $t=\lambda\otimes\omega(f)(x)$. Hal ini mengikuti dengan meninjau
	bayangan $\lambda\otimes x\in
	\omega(Y)^\vee\dotimes{\Bbbk}\omega(A)$ dalam diagram komutatif
	% segment-id: o014.li2.chapter9.reconstruction-a.q007
	\[\begin{tikzcd}
		\omega(Y)^\vee \dotimes{\Bbbk} \omega(A) \arrow[r, "{\omega(f)^\vee \otimes \identity}" inner sep=0.6em] \arrow[d, "{\identity \otimes \omega(f)}"'] & \omega(A)^\vee \dotimes{\Bbbk} \omega(A) \arrow[d] \\
		\omega(Y)^\vee \dotimes{\Bbbk}  \omega(Y) \arrow[r] & \coEnd(\omega).
	\end{tikzcd}\]
\end{proof}

% segment-id: o014.li2.chapter9.reconstruction-a.p006
Untuk kategori kecil secara esensial $\mathcal{A}$ dan funktor
$\omega:\mathcal{A}\to\catesubd{Mod}{\mathrm{pfg}}B$ sembarang,
Definisi--Proposisi~\ref{def:cogebra-L} memfaktorkan $\omega$ secara
kanonik sebagai
% segment-id: o014.li2.chapter9.reconstruction-a.q008
\begin{equation}\label{eqn:L-canonical-decomp}
	\mathcal{A} \xrightarrow{\overline{\omega}} \catesubd{Comod}{\mathrm{pf}}\coEnd(\omega) \xrightarrow{U} \catesubd{Mod}{\mathrm{pfg}}B,
\end{equation}
dengan $U$ funktor lupa. Kita hendak menyelidiki kapan
$\overline\omega$ merupakan ekuivalensi.

% segment-id: o014.li2.chapter9.reconstruction-a.p007
Lema berikut mendasari berbagai teorema rekonstruksi di bawah. Inti
buktinya memakai Teorema Monadisitas Beck~\ref{prop:Beck}. Kita juga
memerlukan pencirian funktor setia eksak dari
Proposisi~\ref{prop:faithful-exact}, konsep kategori abelian hingga secara lokal dari
Definisi~\ref{def:locally-finite}, dan beberapa teknik Ind-isasi dari
\CHref{sec:app-Ind}.

% segment-id: o014.li2.chapter9.reconstruction-a.l001
\begin{lemma}\label{prop:Tannaka-projective-generator}
	Misalkan $\Bbbk$ medan, $\mathcal{A}$ kategori abelian hingga secara lokal, dan
	$s$ generator projektif bagi $\mathcal{A}$. Tinjau funktor setia eksak
	% segment-id: o014.li2.chapter9.reconstruction-a.q009
	\[ \omega: \mathcal{A} \to \cated{Mod}B, \]
	dan andaikan ia mengambil nilai dalam
	$\catesubd{Mod}{\mathrm{pfg}}B$.
	\begin{enumerate}[(i)]
		\item Koaljabar $\coEnd(\omega)$ dapat ditulis secara konkret sebagai
		$P^\vee\dotimes{A}P$, dengan
		$A=\End_{\mathcal{A}}(s)$ dan $P=\omega(s)$.
		\item Funktor $\overline\omega$ dalam faktorisasi kanonik
		\eqref{eqn:L-canonical-decomp} merupakan ekuivalensi kategori.
		\item Berlaku
		$\catesubd{Comod}{\mathrm{pf}}\coEnd(\omega)=
		\catesubd{Comod}{\mathrm f}\coEnd(\omega)$.
	\end{enumerate}
\end{lemma}
\begin{proof}
	Perhatikan bahwa $A$ aljabar-$\Bbbk$ berdimensi hingga. Semua hipotesis
	Teorema~\ref{prop:identify-ModR-fg} terpenuhi, dan
	$\Hom_{\mathcal{A}}(s,\mathord\cdot)$ memberikan ekuivalensi
	$\mathcal{A}\to\catesubd{Mod}{\mathrm{fg}}A$. Jadi masalah direduksi
	ke kasus $\mathcal{A}=\catesubd{Mod}{\mathrm{fg}}A$ dan $s=A$.
	Pengamatan pertama ialah bahwa $\omega$ kini meluas secara kanonik
	menjadi funktor
	% segment-id: o014.li2.chapter9.reconstruction-a.q010
	\[ \Omega: \cated{Mod}A \to \cated{Mod}B, \]
	dengan konstruksi berikut.
	\begin{description}
		\item[Pada objek] Ambil kolimit terfilter kecil
		% segment-id: o014.li2.chapter9.reconstruction-a.q011
		\begin{equation}\label{eqn:Tannaka-tilde-omega}
			\Omega(N) = \varinjlim_{\substack{N' \subset N \\ \text{terbangkit hingga}}} \omega(N').
		\end{equation}
		\item[Pada morfisme] Untuk homomorfisme modul-$A$ kanan
		$f:N_1\to N_2$, ambil
		% segment-id: o014.li2.chapter9.reconstruction-a.q012
		\begin{align*}
			\Omega(f) & = \varprojlim_{N'_1 \subset N_1} \varinjlim_{\substack{N'_2 \subset N_2 \\ N'_2 \supset \Image(f|_{N'_1}) }} \left[ \omega(f|_{N'_1}): \omega(N'_1) \to \omega(N'_2) \right] \\
			& \in \varprojlim_{N'_1 \subset N_1} \varinjlim_{N'_2 \subset N_2} \Hom_{\cated{Mod}B}\left( \omega(N'_1), \omega(N'_2)\right).
		\end{align*}
		Ruas kanan menentukan homomorfisme
		$\Omega(N_1)\to\Omega(N_2)$ dengan cara yang jelas.
	\end{description}

	Jika pembaca menggunakan teori Ind-isasi dari
	\S\S\ref{sec:ind-pro}--\ref{sec:Indization}, konstruksi ini hanya
	mengidentifikasikan $\cated{Mod}A$ dengan
	$\IndC(\catesubd{Mod}{\mathrm{fg}}A)$, lalu memakai kekokompletan
	$\cated{Mod}B$ untuk memperluas $\omega$ ke
	$\IndC(\catesubd{Mod}{\mathrm{fg}}A)$ melalui kolimit terfilter; lihat
	Definisi--Proposisi~\ref{def:Ind-hat-functor}.

	Sifat-sifat berikut merupakan fakta umum tentang Ind-isasi, dan dapat
	pula diperiksa langsung dari konstruksi $\Omega$.
	\begin{itemize}
		\item Pembatasan $\Omega$ pada
		$\catesubd{Mod}{\mathrm{fg}}A$ isomorfik dengan $\omega$.
		\item $\Omega$ merupakan funktor $\Bbbk$-linear eksak; lihat
		Proposisi~\ref{prop:Ind-ext-exactness}(ii) dan pembahasan sebelumnya,
		yang menggunakan keeksakan $\omega$ dan fakta bahwa
		$\cated{Mod}B$ kategori Grothendieck.
		\item $\Omega$ mempertahankan semua kolimit kecil; lihat
		Proposisi~\ref{prop:Ind-ext-exactness}(iii), yang menggunakan fakta
		bahwa $\catesubd{Mod}{\mathrm{fg}}A$ kecil secara esensial.
	\end{itemize}

	Selain itu, karena $\omega$ eksak kiri, untuk semua submodul terbangkit
	hingga $N'\subset N''$ dari $N$, homomorfisme transisi
	$\omega(N')\to\omega(N'')$ dalam \eqref{eqn:Tannaka-tilde-omega}
	bersifat monik. Konstruksi elementer kolimit terfilter dalam kategori
	modul kemudian memberikan
	% segment-id: o014.li2.chapter9.reconstruction-a.q013
	\[ \Omega(N) = 0 \iff \forall N', \; \omega(N') = 0 \stackrel{\because \;\omega\;\text{setia}}{\iff} \; \forall N', \; N' = 0 \iff N = 0. \]
	Jadi $\Omega$ setia dan eksak.

	Sekarang Teorema~\ref{prop:Morita} dapat diterapkan pada $\Omega$ untuk
	memperoleh
	% segment-id: o014.li2.chapter9.reconstruction-a.q014
	\[ \Omega \simeq (\cdot) \dotimes{A} P, \quad P := \Omega(A) = \omega(A) \in \Obj(\catesubd{Mod}{\mathrm{pfg}}B). \]
	Pembatasannya pada $\catesubd{Mod}{\mathrm{fg}}A$ memberikan funktor
	asal $\omega$. Bersama Proposisi~\ref{prop:endomorphism-vs-P}, ini juga
	memberikan deskripsi $\coEnd(\omega)$ dalam (i).

	Karena $\Omega$ setia eksak, Proposisi~\ref{prop:Morita-descent}
	menunjukkan bahwa pasangan adjoin
	% segment-id: o014.li2.chapter9.reconstruction-a.q015
	\begin{equation*}\begin{tikzcd}
		(\cdot) \dotimes{A} P : \cated{Mod}A \arrow[shift left, r] & \cated{Mod}B: (\cdot) \dotimes{B} P^\vee \arrow[shift left, l]
	\end{tikzcd}\end{equation*}
	bersifat komonadik, dengan koaljabar yang bersesuaian
	$P^\vee\dotimes{A}P$. Dengan demikian, $\Omega$ berfaktor secara kanonik
	sebagai ekuivalensi diikuti funktor lupa:
	% segment-id: o014.li2.chapter9.reconstruction-a.q016
	\[ \cated{Mod}A \xrightarrow[\sim]{\overline{\Omega}} \cated{Comod}\left(P^\vee \dotimes{A} P\right) \xrightarrow{U} \cated{Mod}B. \]

	Menurut Proposisi~\ref{prop:Morita-comonad-ft}, faktorisasi ini membatasi
	menjadi diagram komutatif
	% segment-id: o014.li2.chapter9.reconstruction-a.q017
	\[\begin{tikzcd}
		\catesubd{Mod}{\mathrm{fg}}A \arrow[r, "\sim"'] \arrow[rd, bend right] & \catesubd{Comod}{\mathrm{f}} \left(P^\vee \dotimes{A} P\right) \arrow[r, "U"]  \arrow[d, "\sim" sloped, "{\text{(i)}}"'] & \catesubd{Mod}{\mathrm{fg}}B. \\
		& \catesubd{Comod}{\mathrm{f}}\coEnd(\omega) &
	\end{tikzcd}\]
	Komposisi baris pertama ialah $\omega$, yang diketahui mengambil nilai
	dalam $\catesubd{Mod}{\mathrm{pfg}}B$; ini membuktikan (iii). Panah
	melengkung dalam diagram tersebut kemudian sama dengan
	$\overline\omega$, sehingga (ii) mengikuti.
\end{proof}
